\section{Experiment support: data and metrics}
\label{sec:platform-data}

Goal (iv) of the pre-study was the data generator and the metric functions
that Experiment~2 needs. Both are part of the platform rather than of the
experiment script, so that any method registered in the harness can be
trained on the same data and scored by the same functions. This section
describes the mechanisms; the values chosen for Experiment~2 (problem size,
$\lambda$, the three widths, dataset sizes) are given with their rationale in
Chapter~\ref{ch:exp2}.

\subsection{LASSO dataset generator}

\texttt{Benchmarking/data/lasso.py} generates, for one seed, a family of
datasets that share a single dictionary $A \in \mathbb{R}^{m \times n}$ with
i.i.d.\ $\mathcal{N}(0, 1/m)$ entries and unit-norm columns. Sparse signals
have exactly $\mathrm{round}(p\,n)$ nonzero coordinates on a uniformly random
support with $\mathcal{N}(0,1)$ values (a Bernoulli-support variant matching
the Primer's own convention is also provided), and measurements
$b = A x_\mathrm{true} + \varepsilon$ with $\varepsilon$ scaled to a requested
SNR, or omitted for the noiseless case. The sparsity level $p$ is the quantity
whose distribution the experiment controls: a training set is generated from a
\emph{sparsity specification} that is either a fixed $p$ or an interval from
which $p$ is drawn uniformly \emph{per instance}, and the generator writes one
training and validation set per specification (the ``narrow'', ``medium'' and
``wide'' widths of Experiment~2) plus a battery of test sets at fixed $p$
values. Because $A$ and the test battery are shared across widths, training
width is the only thing that differs between the runs that are later
compared.

Two further files are written per split. The seeded starting point
\texttt{*\_x0.npy} of Section~\ref{sec:platform-repro} is generated here, once,
so that it is a property of the data rather than of any method. For ALISTA,
whose weights are analytic rather than learned, the generator computes the
matrix $W$ from $A$ by solving the coherence-minimisation problem of Liu et
al.~\cite{liu2019alista} exactly, one linear program per column
($w_i = \arg\min_w \max_{j \neq i} |w^\top a_j|$ subject to $w^\top a_i = 1$),
and stores it next to $A$.

\subsection{Metric functions and the reference solution}

\texttt{core/lasso\_metrics.py} implements the four metrics of the
pre-study (Section~\ref{sec:prestudy-exp2}) for a batch of instances and a
batch of predicted solutions $x$:
the suboptimality gap $g_q(x) = f_q(x) - f_q^\ast$
(Eq.~\ref{eq:subopt-gap}); the modified relative loss
$R_{f} = \mathbb{E}_q[f_q(x) - f_q^\ast] / \mathbb{E}_q[f_q^\ast]$
(Eq.~\ref{eq:rel-loss}), a ratio of expectations over the test set as in the
Primer; the NMSE against $x_\mathrm{true}$ in dB (Eq.~\ref{eq:nmse}); and the
LASSO-optimal recovery error, the same NMSE evaluated at the LASSO minimiser
$x^\ast$ itself, which measures how well the objective at this $\lambda$ can
recover the signal at all and thereby separates a hard test sparsity from a
weak method. The three per-instance quantities are returned as a mean and a
standard deviation over the test set; $R_f$, being a ratio of two means, is a
single number.

Three of the four depend on the exact minimiser $x^\ast$ and its objective
value $f^\ast$, so the quality of the reference solver bounds the precision
of every metric. The solver is a batched FISTA in double precision with
per-instance adaptive restart~\cite{odonoghue2015restart} and a
\emph{certified} stopping rule: every 25 iterations it evaluates the LASSO
duality gap from the current iterate's dual point, which upper-bounds
$f(x) - f^\ast$ for any $x$, and stops when the worst instance in the batch
is within a relative gap of $10^{-10}$, up to an iteration cap. A predicted
$x$ that beats the certified $x^\ast$ by more than numerical tolerance is
therefore impossible, and the function warns if it happens, which catches an
$x$ scored under the wrong $\lambda$. Because $x^\ast$ depends only on
$(A, b, \lambda)$ and the tolerances, it is computed once and cached under a
key that is the SHA-256 hash of the actual $A$ and $b$ bytes, not of a
filename, so that a cache entry can never be reused for a different batch by
mistake. All methods, widths and seeds that evaluate the same test file share
the same $x^\ast$.

\subsection{Classical baselines and per-iteration curves}

Two analysis-time tools complete the experiment support. First,
\texttt{classical\_baselines.py} runs textbook ISTA and FISTA (fixed step
$1/L$, no restart, from the shared $x_0$) on exactly the test instances and
seeds of a sweep and scores them with the same metric functions and the same
cached $x^\ast$, so that a learned optimizer's row can be placed next to a
classical row produced under identical conditions. It records the metrics at
chosen iteration counts (16, 100 and 1000 in Chapter~\ref{ch:exp2}) and the
full mean objective trajectory. Second, the model-free evaluation scripts
save the batch-mean objective at every step, and
\texttt{plot\_compute\_matched\_rf.py} turns that record into
$R_f$ as a function of the iteration $k$, with the unrolled networks' values
at their fixed depth as reference lines. Together these give the
matched-budget view used in Chapter~\ref{ch:exp2}, in which iteration $k$
means $k$ learned steps, $k$ solver iterations or $k$ unrolled layers, all
from the same $x_0$; the caveat that equal iteration counts are not equal
compute is stated there.
